\section{IRQ 分发框架：\texttt{irq.c}}

\subsection{函数指针数组}

\codefile{src/kernel/arch/irq.c}
\begin{lstlisting}[style=cstyle]
static irq_handler_t g_irq_handlers[16];  // 初始全为 NULL

void irq_install_handler(uint8_t irq, irq_handler_t handler) {
    if (irq < 16) {
        g_irq_handlers[irq] = handler;
    }
}
\end{lstlisting}

\begin{figure}[H]
  \centering
  % figures/ch05/fig03-irq-g_irq_handlers.tex — auto-extracted from chapter source
\begin{tikzpicture}[node distance=0cm]
  % Array
  \foreach \i in {0,...,7} {
    \node[memcell, minimum width=2.5cm, minimum height=0.5cm] (h\i) at (0, -\i*0.5) {};
    \node[memaddr, left=0.1cm of h\i, font=\tiny\ttfamily] {[\i]};
  }
  \node[font=\small\ttfamily, above=0.2cm of h0] {g\_irq\_handlers[16]};

  % NULL for most
  \foreach \i in {0,2,3,4,5,6,7} {
    \node[font=\tiny\ttfamily, text=gray] at (h\i) {NULL};
  }

  % Handler for IRQ1
  \node[font=\tiny\ttfamily, text=blue!60!black] at (h1) {keyboard\_irq1\_handler};

  % Arrow to handler function
  \node[draw, minimum width=3.5cm, fill=green!10, right=3cm of h1, font=\small] (kbdfn) {keyboard\_irq1\_handler()};
  \draw[pointer] (h1.east) -- (kbdfn.west);
\end{tikzpicture}

  \caption{IRQ 分发表：\texttt{g\_irq\_handlers} 函数指针数组}
\end{figure}

\subsection{统一 C 入口}

\codefile{src/kernel/arch/irq.c}
\begin{lstlisting}[style=cstyle]
void irq_handler_c(const irq_regs_t* r) {
    uint8_t irq = (uint8_t)(r->int_no - 0x20u);

    // 调用已注册的 handler（如果有）
    if (irq < 16 && g_irq_handlers[irq]) {
        g_irq_handlers[irq](r);
    }

    // 统一发送 EOI
    pic_send_eoi(irq);
}
\end{lstlisting}

\begin{designbox}
EOI 在 dispatcher 中统一发送，而非让每个 handler 自己发。好处：
\begin{enumerate}
  \item 驱动开发者不会忘记发 EOI
  \item handler 只需专注于设备逻辑
  \item 将来做更复杂的 IRQ 优先级管理时只改 dispatcher 一处
\end{enumerate}
\end{designbox}

\subsection{IRQ 汇编 Stub}

IRQ stub 和 ISR stub 结构相同（push err\_code=0 → push int\_no → pushad → call C）：

\codefile{src/kernel/arch/irq\_stubs.asm}
\begin{lstlisting}[style=nasmstyle]
%macro IRQ 2
irq%1:
    push dword 0          ; err_code = 0（硬件 IRQ 无 error code）
    push dword %2         ; int_no = vector (0x20 + irq)
    jmp irq_common_stub
%endmacro

IRQ 0, 0x20    ; IRQ0 → 向量 0x20
IRQ 1, 0x21    ; IRQ1 → 向量 0x21（键盘）
; ... 共 16 个 ...
\end{lstlisting}

\texttt{irq\_common\_stub} 和 \texttt{isr\_common\_stub} 完全一样
（pushad → push 段寄存器 → call C → 恢复 → iret），
只是调用的 C 函数不同（\texttt{irq\_handler\_c} vs \texttt{isr\_handler\_c}）。

% ============================================================================

